\documentclass[final, english, english, a4paper, 12pt, numbers=noenddot, % Bei Änderung der Sprache 2 x kompilieren!
% tudscr spezifische Einstellungen
cd=false,
cdfont=false,cdfont=nohead,
cdmath=false,
cdhead=false,
cdfoot=true,
cdcover=monochrome,
cdgeometry=asymmetric,
declaration=heading,
declaration=notoc,
abstract=heading,
]{tudscrreprt}

%################ Hilfreiche Pakete laden #######################

%################ Hilfreiche Pakete laden #######################

\usepackage{settings/tudbwlimPackages}
\usepackage{settings/tudbwlimStyle}
\usepackage{scrhack}
\usepackage{microtype}
\usepackage{tikz,pgfplots}
\pgfplotsset{compat=newest}
\usepackage[chapter]{algorithm} % Falls das Paket floatrow geladen wird, muss dieses Paket danach geladen werden.
\iflanguage{english}{\floatname{algorithm}{Algorithm}\renewcommand{\listalgorithmname}{List of Algorithms}}{\floatname{algorithm}{Algorithmus}\renewcommand{\listalgorithmname}{Algorithmenverzeichnis}} % Algorithm-Umgebung an die verwendete Sprache anpassen
\usepackage{comment}
\usepackage{subcaption}
\usepackage{caption}
\usetikzlibrary{positioning}
\usetikzlibrary{mindmap}
\usepackage{graphicx}
\usepackage{float}
\usepackage{longtable}
\usepackage{xparse} % for \NewDocumentCommand
\usepackage{multirow}
\usepackage{array}
\usepackage{makecell}
\usepackage[justification=centering]{caption} % Für die Zentrierung der Bildunterschrift

%################ Notwendige Pakete laden #######################


\usepackage[bibencoding=auto,citestyle=authoryear-ibid,bibstyle=authoryear,maxcitenames=3,maxbibnames=10, backend = biber]{biblatex}% Vergessen Sie nicht in den Optionen das Bibliographieprogramm auf "biber" umzustellen! Um die Vorlage mit BibTeX nutzen zu können, muss die Option "backend=bibtex" übergeben werden. Es ist jedoch biber zu empfehlen, beachten Sie dazu die Hinweise der biblatex-Paketdokumentation im Abschnitt 3.15 "Using the fallback BibtTeX backend".
\usepackage{settings/BiblatexSetup}

\AfterPackage*{biblatex}%
{
    \RequirePackage[breaklinks=true, colorlinks=true, linktoc=section, linkcolor=blue, citecolor=black, hidelinks]{hyperref}
    % Da hyperref allerhand Veränderungen an vielen Standardbefehlen vornimmt, sollte dieses als letztes in der Präambel eingebunden werden. Nur Pakete, bei denen in der Dokumentation explizit darauf hingewiesen wird, dass diese nach hyperref zu laden sind, sollten auch danach folgen.
    \hypersetup{pdfprintscaling=None} % gleiches Verhalten, auch ohne hyperref, liefert: \pdfcatalog{/ViewerPreferences<</PrintScaling/None>>}
    \usepackage{footnote} % https://tex.stackexchange.com/questions/207192/footcite-in-float-caption
    \makesavenoteenv{figure}
    \makesavenoteenv{table}
    \makesavenoteenv{algorithm}
}


\AfterPackage*{hyperref}
{
    \RequirePackage[automake,acronym,symbols,nomain,translate=babel]{glossaries}
    \usepackage{settings/GlossariesSetup}
}
%############## Own commands #####################´
\newcommand{\parbreak}{\vspace{\baselineskip}\noindent}
\newcommand{\gendreauDataSet}{Gendreau (2006)}
\newcommand{\mouraDataSet}{Moura (2009)}
\newcommand{\ceschiaDataSet}{Ceschia (2013)}
\newcommand{\krebsADataSet}{Krebs (2021a)}
\newcommand{\krebsBDataSet}{Krebs (2021b)}
\newcommand{\gendreauDataSetText}{\gendreauDataSet\ }
\newcommand{\mouraDataSetText}{\mouraDataSet\ }
\newcommand{\ceschiaDataSetText}{\ceschiaDataSet\ }
\newcommand{\krebsADataSetText}{\krebsADataSet\ }
\newcommand{\krebsBDataSetText}{\krebsBDataSet\ }

%############## Ideas for packages #####################´
% Folgende auskommentierte Pakete sind als Vorschläge zu verstehen. Für Funktionsweise und Anwendungsfälle wird auf das Benutzerhandbuch "tudscr.pdf" (http://mirrors.ctan.org/macros/latex/contrib/tudscr/doc/tudscr.pdf) oder auf die entsprechende CTAN Dokumentation verwiesen.
%\usepackage{listings} %For Visualization of programming source code 
%\lstset{%
%	inputencoding=utf8,extendedchars=true,
%	literate=%
%	{ä}{{\"a}}1 {ö}{{\"o}}1 {ü}{{\"u}}1
%	{Ä}{{\"A}}1 {Ö}{{\"O}}1 {Ü}{{\"U}}1
%	{~}{{\textasciitilde}}1 {ß}{{\ss}}1
%	}

%\usepackage{calc} % Can be used to calculate inline distances in .tex files

%################ Sonstige Einstellungen/Befehle #####################
\onehalfspacing

%################ Abkürzungen #######################
\glsenableentrycount % aktiviert \cgls, \cglspl, \cGls, \cGlspl, siehe https://tex.stackexchange.com/questions/98494/glossaries-dont-print-single-occurences/230664#230664

% Abkürzungen, die im Abkürzungsverzeichnis auftauchen und automatisch durch das glossaries-Paket sortiert werden
\newacronym[description={Container loading problem}]{CLP}{CLP}{container loading problem}
\newacronym[description={Vehicle routing problem}]{VRP}{VRP}{vehicle routing problem}
\newacronym[description={Capacitated vehicle routing problem}]{CVRP}{CVRP}{capacitated vehicle routing problem}
\newacronym[description={Bin packing problem}]{BPP}{BPP}{bin packing problem}
\newacronym[description={Last-in-first-out}]{LIFO}{LIFO}{last-in-first-out}
\newacronym[description={Cutting and packing}]{CaP}{C\&P}{cutting and packing}
\newacronym[description={Constraint programming}]{CP}{CP}{constraint programming}
\newacronym[description={Mixed-integer programming}]{MIP}{MIP}{mixed-integer programming}
\newacronym[description={Genetic algorithm}]{GA}{GA}{genetic algorithm}
\newacronym[description={Tabu search}]{TS}{TS}{tabu search}
\newacronym[description={Simulated annealing}]{SA}{SA}{simulated annealing}
\newacronym[description={Greedy randomized adaptive search procedure}]{GRASP}{GRASP}{greedy randomized adaptive search procedure}
\newacronym[description={Variable neighborhood search}]{VNS}{VNS}{variable neighborhood search}
\newacronym[description={Pallet loading problem}]{PLP}{PLP}{pallet loading problem}
\newacronym[description={Machine learning}]{ML}{ML}{machine learning}
\newacronym[description={Artificial neural network}]{ANN}{ANN}{artificial neural network}
%\newacronym[description={Support vector machine}]{SVM}{SVM}{support vector machine}
\newacronym[description={Logistic regression}]{LR}{LR}{logistic regression}
\newacronym[description={Two-dimensional loading capacitated vehicle routing problem}]{2L-CVRP}{2L--CVRP}{two-dimensional loading capacitated vehicle routing problem}
\newacronym[description={Three-dimensional loading capacitated vehicle routing problem}]{3L-CVRP}{3L--CVRP}{three-dimensional loading capacitated vehicle routing problem}
\newacronym[description={Three-dimensional loading vehicle routing problem with time windows}]{3L-VRPTW}{3L--VRPTW}{three-dimensional loading vehicle routing problem with time windows}
\newacronym[description={Time windows}]{TW}{TW}{time windows}
%\newacronym[description={Lower bound}]{LB}{LB}{lower bound}
\newacronym[description={Upper bound}]{UB}{UB}{upper bound}
\newacronym[description={Manual last-in-first-out}]{MLIFO}{MLIFO}{manual last-in-first-out}
\newacronym[description={Load-bearing strength}]{LBS}{LBS}{load-bearing strength}


% Abkürzungen, die nicht im Abkürzungsverzeichnis aufgeführt werden
%\newabbreviation{zB}{z.\,B.}{zum Beispiel}

% Do i need this? 
\glsaddall[types=abbreviation]

% Symbole die im Symbolverzeichniss erscheinen sollen
% Mit "F5" kompilieren oder "Tools -> Befehle -> makeglossaries" (F9) starten, um das Symbolverzeichnis zu aktualisieren
% \newsymb{<sort by>}{<name>}{<symbol>}{<unit>}
%
\begin{comment}
\newsymb{E}{Erwartungswert}{\mathbb{E}\left(\cdot\right)}{}
\newsymb{P}{Wahrscheinlichkeitsmaß}{\mathbb{P}\left(\cdot\right)}{}
\newsymb{V}{Varianz}{\mathbb{V}\left(\cdot\right)}{}
\newsymb{X}{Zufallsvariable}{X}{}
%
% Oder so verwenden und im Fließtext dann mit $\ExpValue$ arbeiten:
%\newcommand{\ExpValue}{\mathbb{E}\left(\cdot\right)}
%\newsymb{E}{Erwartungswert}{\ExpValue}{}
%
\glsaddall[types={symbols}] % Alle Symbole werden dem Symbolverzeichnis hinzugefügt
\end{comment}

%################ Bibliographie laden #######################
\addbibresource{./References.bib} % Pfad/Name der .bib-Datei

%################ Ende Präambel #######################
\pagenumbering{Roman}

\begin{document}
%################ Aufruf Deckblatt #######################
% mögliche Optionen, für weitere Informationen, siehe S. 23 des Benutzerhandbuchs des tudscr-Paket (http://mirrors.ctan.org/macros/latex/contrib/tudscr/doc/tudscr.pdf):

%%%%%%%%%%%%%%%%%%%%%%%%%%%%%%%%%%%%%
%\thesis{\diplomathesisname}		% Diplomarbeit/Diploma-Thesis
%\graduation[Dipl.-Kffr.]{Diplom-Kauffrau} % [Dipl.-Kfm.]{Diplom-Kaufmann}
%
%\thesis{\masterthesisname}			% Master-Arbeit/Master Thesis
%\graduation[M.Sc.]{Master of Science}
%
%\thesis{\bachelorthesisname}		% Bachelor-Arbeit/Bachelor Thesis
%\graduation[B.Sc.]{Bachelor of Science}

\renewcommand{\seminarCategoryText}{As part of the Industrial Management research seminar}
%
\newsavebox\thesisbox
\savebox\thesisbox{\normalfont\large\sffamily \parbox{\textwidth}{\vspace{6pt} \raggedright \seminarCategoryText}}
% https://github.com/tud-cd/tudscr/issues/85 %
%
\thesis{\seminarpapername \seminarCategory{\\\usebox\thesisbox}}		% Seminararbeit/Seminar Paper
%%%%%%%%%%%%%%%%%%%%%%%%%%%%%%%%%%%%%

\title{Feasibility prediction for the container loading problem}
%\subtitle{Optionaler Unter-/Zusatztitel}
\author{%
	Maximilian Hubmann
	\matriculationnumber{4805234}
	\dateofbirth{21.03.2000}
	\placeofbirth{Nürnberg}
	\course{Wirtschaftsingenieurwesen, 9. FS}
}
\date[]{\today}
\supervisor{Dipl.-Wi.-Ing. Florian Linß}
\professor{Prof. Dr. Udo Buscher}

\setcounter{page}{1}

%%% IM 
\headlogo{./pictures/IM-Logo}
\faculty{Faculty of business and economics}
\chair{Chair of Business Administration, in particular Industrial Management}
\maketitle[cdtitle=true, cdfont=false]

%################ Abstract ######################
%\begin{abstract}
%Das ist ja maal so cool, dass ich ich einfahc hier einen Tsext schreiben kann
%\end{abstract}

%%% Danksagung
%\danke{Danksagung}{Text}

%####################### Verzeichnisse ################%
%\microtypesetup{protrusion=false}
% Inhaltsverzeichnis
\tableofcontents

% Abbildungsverzeichnis (falls nichts benötigt, einfach als Kommentar setzen)
%\listoffigures
%\addcontentsline{toc}{chapter}{\listfigurename}

% Tabellenverzeichnis (falls nichts benötigt, einfach als Kommentar setzen)
%\listoftables
%\addcontentsline{toc}{chapter}{\listtablename}

% Algorithmenverzeichnis (falls nichts benötigt, einfach als Kommentar setzen)
%\listofalgorithms
%\addcontentsline{toc}{chapter}{\listalgorithmname}

\microtypesetup{protrusion=true}

%Abkürzungsverzeichnis (falls nichts benötigt, einfach als Blockkommentar setzen)
%\printacronyms[style=bwlimsuper]
%\addcontentsline{toc}{chapter}{\acronymname}

% Symbolverzeichnis (falls nichts benötigt, einfach als Blockkommentar setzen)
%\printsymbols[style=symblong, title=\listsymbolname]
%\addcontentsline{toc}{chapter}{\listsymbolname}


%############### Einstellungen für Fließtext setzen ####################
\clearpage
\setcounter{page}{1}\pagenumbering{arabic}


%####################### Fließtext ####################

% Introduction
\chapter{Erstes Kapitel}

Das \cgls{scm} eines \glsuseri{kmu} ist auszubauen. \cGls{mw} ist groß. \cGls{mw} sollte \acrshort{zB} am Anfang eines Satzes groß geschrieben werden. \glqq \Acrlong{zB}\grqq\xspace kann auch ausgeschrieben werden. Die Befehle \verb|\cgls{}| bzw. \verb|\cGls{}| zählen die Verwendung einer Abkürzung und fügen sie dem Abkürzungsverzeichnis hinzu, wenn sie mindestens zweimal verwendet wird. Text Text Text Text Text Text Text Text Text Text Text Text Text Text Text Text Text Text Text Text Text Text Text Text Text Text Text Text Text Text Text Text Text Text Text Text Text Text Text Text Text Text Text Text Text Text Text Text Text Text \glssymbol{E} Text \glssymbol{V} Text \glssymbol{P} Text \glssymbol{X} Text Text Text Text Text Text Text Text Text Text Text Text.

Text \cgls{crm} Text Text Text Text Text Text Text Text Text Text Text Text Text Text Text Text Text Text Text Text Text Text Text Text Text Text Text Text Text Text Text Text Text Text Text Text Text Text Text Text Text Text Text Text Text Text Text Text Text Text Text Text Text Text Text Text Text Text Text Text Text Text Text Text Text Text Text Text Text Text Text Text Text Text Text Text Text Text Text Text Text Text Text Text Text Text Text Text Text Text Text Text Text Text Text Text Text Abbildung~\ref{fig:1}.%

\begin{figure}[ht]
    \centering
    \includegraphics[width=0.1\textwidth]{./pictures/IM-Logo}
    \caption[Erste Abbildung]{Erste Abbildung\footnotemark}\label{fig:1}
\end{figure}\footnotetext{Eigene Abbildung.}

Text Text Text Text Text Text Text Text Text Text Text Text Text Text Text Text Text Text Text Text Text Text Text Text Text Text Text Text Text Text Text Text Text Text Text Text Text Text Text Text Text Text Text Text Text Text Text Text Text Text Text Text Text Text Text Text Text Text Text Text Text Text Text Text Text Text Text Text Text Text Text Text Text Text Text Text Text Text Text Text Text Text Text Text Text Text Text Text Text Text Text Text Text Text Text Text Text Text Text Text Text Text Text Text Text Text Text Text Text Text Text \cgls{scm} Text Text Text Text Text Text Text Text Text Text Text Text Text Text Text Text Text Text Text Text Text Text Text Text Text Text Text Text Text Text Text Text Text Text Text Text Text Text Text Text Text Text Text Text Text Text Text Text Text Text Text Text Text Text Text Text Text Text Text Text Text Text Text Text\footnote{Fußnotentext.} Text Text Text Text Text Text Text Text Tabelle~\ref{tab:1}.%

\begin{table}[tb]
    \centering
    \begin{tabular}{ccc}
        \toprule
        A & B & C \\
        \midrule
        1 & 2 & 3 \\
        4 & 5 & 6 \\
        \bottomrule
    \end{tabular}
    \caption{Erste Tabelle}\label{tab:1}
\end{table}
Text Text Text Text Text Text Text Text Text Text Text Text Text Text Text Text Text Text Text Text Text Text Text Text Text Text Text Text Text Text Text Text Text Text Text Text Text Text Text Text Text Text Text \cgls{kmu} Text Text Text Text Text Text Text Text Text Text Text Text Text Text Text Text Text Text Text Text Text Text Text Text Text Text Text Text Text Text Text Text Text Text Text Text Text Text Text Text Text Text Text Text Text Text Text Text Text Text Text Text. Wie \citeauthor{HinzH1:2009} zeigen, Text Text.\footcite[Vgl.][33]{HinzH1:2009} Text Text \cgls{crm}.

\begin{figure}[h]
    \centering
    \begin{tikzpicture}%[scale=0.975]
        \begin{axis}[
                xmin=-0.5, xmax=5,
                ymin=-0.5, ymax=5,
                xtick={1,...,4},
                ytick={1,...,4},
                axis y line=center,
                axis x line=middle,
                ylabel={$y$},
                xlabel={$x$},
                xlabel style={below right},
                ylabel style={above left},
            ]

            \node (A) at (axis cs:2,1) {};
            \node[right] at (A) {$A$};
            \draw[fill] (A) circle [radius=1.75pt];

        \end{axis}
    \end{tikzpicture}
    \caption[Abbildungsunterschrift]{Abbildungsunterschrift\footcite[Vgl.][39]{HinzH1:2009}}
\end{figure}

Text Text Text Text Text Text Text Text Text Text Text Text Text Text Text Text Text Text Text Text Text Text Text Text Text Text Text Text Text Text Text Text Text Text Text Text Text Text Text Text Text Text Text Text Text Text Text Text Text Text Text Text Text Text Text Text Text Text Text Text Text Text Text Text Text Text Text Text Text Text Text Text Text Text Text Text Text Text Text Text Text Text Text Text Text Text Text Text Text Text Text Text Text. Verweis mit mehreren Quellenangaben.\footcites(Vgl.)()[][24]{HinzH1:2009}[][15]{KunzK1:2010}



% Main Part
\chapter{Zweites Kapitel}

\section{Erster Abschnitt des zweiten Kapitels}

Text Text Text Text Text Text Text Text Text Text Text Text Text Text Text Text Text Text Text Text Text Text Text Text Text Text Text Text Text Text Text Text Text Text Text Text Text Text Text Text Text Text Text Text Text Text Text Text Text Text Text Text Text Text Text Text Text Text Text Text Text Text Text Text Text Text Text Text Text Text Text Text Text Text Text Text Text Text Text Text Text Text Text Text Text Text Text Text Text Text Text Text Text Text Text Text Text Text Text Text Text Text Text. Wie im Algorithmus~\ref{alg:euclid} dargestellt\dots

\begin{algorithm}[htb]
	\caption{Euclid's Algorithmus}\label{alg:euclid}
	\begin{algorithmic}[1]
		\Procedure{Euclid}{$a,b$}\Comment{Der größte gemeinsame Teiler (ggT) von $a$ und $b$}
			\State $r \gets a \bmod b$
			\While{$r\not=0$}\Comment{Das Ergebnis liegt vor wenn $r=0$}
				\State $a \gets b$
				\State $b \gets r$
				\State $r \gets a \bmod b$
			\EndWhile
			\State \textbf{return} $b$ \Comment{Der ggT ist $b$}
		\EndProcedure
	\end{algorithmic}
\end{algorithm}

Text Text Text Text Text Text Text Text Text Text Text Text Text Text Text Text Text Text Text Text Text Text Text Text Text Text Text Text Text Text Text Text Text Text Text Text Text Text Text Text Text Text Text Text Text Text Text Text Text Text Text Text Text Text Text Text Text Text Text Text Text Text Text Text Text Text Text Text Text Text Text Text Text Text Text Text Text Text Text Text Text Text Text Text Text Text Text Text Text Text Text Text Text Text Text Text Text Text Text Text Text Text Text.

\section{Zweiter Abschnitt des zweiten Kapitels}

Text Text Text Text Text Text Text Text Text Text Text Text Text Text Text Text Text Text Text Text Text Text Text Text Text Text Text Text Text Text Text Text Text Text Text Text Text Text Text Text Text Text Text Text Text Text Text Text Text Text Text Text Text Text Text Text Text Text Text Text Text Text Text Text Text Text Text Text Text Text Text Text Text Text Text Text Text Text Text Text Text Text Text Text Text Text Text Text Text Text Text Text Text Text Text Text Text Text Text Text Text Text Text.

\section{Dritter Abschnitt des dritten Kapitels}

Text Text Text Text Text Text Text Text Text Text Text Text Text Text Text Text Text Text Text Text Text Text Text Text Text Text Text Text Text Text Text Text Text Text Text Text Text Text Text Text Text Text Text Text Text Text Text Text Text Text Text Text Text Text Text Text Text Text Text Text Text Text Text Text Text Text Text Text Text Text Text Text Text Text Text Text Text Text Text Text Text Text Text Text Text Text Text Text Text Text Text Text Text Text Text Text Text Text Text Text Text Text Text.

Geht das jetzt? Ja :D 


%####################### Appendix #########################
%%############################# Anhang #################################
%\appendix
%\chapter{Mathematischer Anhang}
%\chapter{Programmcodes}

\clearpage
%####################### Literaturverzeichnis #########################
\addcontentsline{toc}{chapter}{\bibname}
\printbibliography
%####################### Selbständigkeitserklärung ####################

\newConfirmationEnglish{Dresden}

%####################### Platzhalter - Abkürzungsverzeichnis muss erstellt werden####################

\printacronyms[style=bwlimsuper]

%####################### End of Document ####################
\end{document}